\subsection{Evaluation Metrics}
\label{sec:evaluation_metrics}

The evaluation of intersectional multiclass classification performance ($K = 6$) at the subgroup level is formulated through a One-vs-Rest (OvR) scheme for each demographic group $c \in \{1, \dots, K\}$~\cite{ref44}. This formulation utilizes the confusion matrix components $TP_c$, $TN_c$, $FP_c$, and $FN_c$~\cite{ref45}, \cite{ref46}. Subgroup Accuracy in \eqref{eq:13} measures the ratio of correct predictions, whereas Precision in \eqref{eq:14}, Recall or True Positive Rate (TPR) in \eqref{eq:15}, and F1-Score in \eqref{eq:16} map detection precision and sensitivity~\cite{ref47}. To evaluate the held-out test data, subgroup metrics are aggregated to the global level. Global Accuracy in \eqref{eq:17} calculates the total ratio of correct predictions, whereas Macro Precision in \eqref{eq:18}, Macro Recall in \eqref{eq:19}, and Macro F1-Score in \eqref{eq:20} are averaged without weights across the $K = 6$ subgroups to ensure evaluation parity.

\begin{equation}
\text{Accuracy}_c = \frac{TP_c + TN_c}{TP_c + TN_c + FP_c + FN_c}
\label{eq:13}
\end{equation}
\begin{equation}
\text{Precision}_c = \frac{TP_c}{TP_c + FP_c}
\label{eq:14}
\end{equation}
\begin{equation}
\text{Recall}_c = \frac{TP_c}{TP_c + FN_c}
\label{eq:15}
\end{equation}
\begin{equation}
\text{F1-Score}_c = \frac{2 \cdot \text{Precision}_c \cdot \text{Recall}_c}{\text{Precision}_c + \text{Recall}_c}
\label{eq:16}
\end{equation}
\begin{equation}
\text{Accuracy}_{\text{Global}} = \frac{\sum_{c=1}^K TP_c}{N}
\label{eq:17}
\end{equation}
\begin{equation}
\text{Precision}_{\text{Macro}} = \frac{1}{K} \sum_{c=1}^K \text{Precision}_c
\label{eq:18}
\end{equation}
\begin{equation}
\text{Recall}_{\text{Macro}} = \frac{1}{K} \sum_{c=1}^K \text{Recall}_c
\label{eq:19}
\end{equation}
\begin{equation}
\text{F1-Score}_{\text{Macro}} = \frac{1}{K} \sum_{c=1}^K \text{F1-Score}_c
\label{eq:20}
\end{equation}

The validation of comparative significance between models is evaluated through McNemar's paired inferential hypothesis test with Edwards continuity correction in \eqref{eq:21}~\cite{ref44}. The $\chi^2$ statistic is computed on a 2 $\times$ 2 contingency table based on discordant pairs, where $b$ and $c$ record differences in correct--incorrect classifications between models, reinforced by a two-sided binomial exact test. Empirical accuracy uncertainty on the held-out test data is measured through Wilson Score 95\% Confidence Interval (CI) estimation in \eqref{eq:22}, with $\hat{p}$ denoting the accuracy proportion and $z = 1.96$ as the standard critical value. The evaluation of performance disparity is quantified through Equal Opportunity Difference (EOD) in \eqref{eq:23} as the maximum absolute difference in TPR between subgroup pairs $i, j \in \{1, \dots, K\}$, which is calibrated strictly as an empirical inter-subgroup disparity indicator rather than a mathematical guarantee of universal fairness.

\begin{equation}
\chi^2 = \frac{(|b - c| - 1)^2}{b + c}
\label{eq:21}
\end{equation}
\begin{equation}
\text{CI}_{95\%} = \frac{\hat{p} + \frac{z^2}{2N} \pm z \sqrt{\frac{\hat{p}(1-\hat{p})}{N} + \frac{z^2}{4N^2}}}{1 + \frac{z^2}{N}}
\label{eq:22}
\end{equation}
\begin{equation}
\text{EOD} = \max_{i, j \in \{1, \dots, K\}} |\text{Recall}_i - \text{Recall}_j| = \max \text{TPR} - \min \text{TPR}
\label{eq:23}
\end{equation}
